% Beispiel Dokument
\documentclass[12pt,digital]{hska}

\addbibresource{document.bib}

% Erstellung eines Index und Abkürzungsverzeichnisses/Glossars aktivieren ------------------------------
% \makeindex{}
% \makenomenclature{}

% Das eigentliche Dokument -----------------------------------------------------------------------------
% ------------------------------------------------------------------------------------------------------
\begin{document}

% Titelblatt
% ------------------------------------------------------------------------------------------------------

% vorgegebene Seite verwenden, zB bei Bachelor- oder Masterarbeit:
% ------------------------------------------------------------------------------------------------------
% \thispagestyle{plain}
% \begin{titlepage}
% \includepdf[pages={1}]{sections/coverpage.pdf}
% \end{titlepage}

% einfache Seite selbst schreiben:
% ------------------------------------------------------------------------------------------------------
\thispagestyle{plain}
\begin{titlepage}
\begin{center}
% Logo der HSKA
\includegraphics[width=0.55\textwidth]{images/hka_logo.pdf}\\[12ex]

\LARGE{\textbf{Seminararbeit}\\Sommersemester 2019}\\[8ex]

\textbf{Titel der Arbeit\\
ggf. zweizeilig}\\[12ex]

\normalsize{}
\begin{tabular}{lll}
Student:  & \quad Max Muster & xxxxxx\\[3ex]
Universität: & \quad Hochschule Karlsruhe – Technik und Wirtschaft &\\[1ex]
Betreuer:  & \quad Prof. Dr. Daniel Düsentrieb &\\[3ex]
In Kooperation mit: & \quad Tolle Firma GmbH &\\[1ex]
Betreuer:  & \quad Donald Duck &\\[1ex]
\end{tabular}

\end{center}
\end{titlepage}

% Seitennummerierung -----------------------------------------------------------------------------------
%   Vor dem Hauptteil werden die Seiten in großen römischen Ziffern nummeriert.
% ------------------------------------------------------------------------------------------------------
\pagenumbering{Roman}
\phantomsection{} % Sorgt für korrekte Aufnahme des Inhaltsverzeichnisses in das Inhaltsverzeichnis
\addcontentsline{toc}{section}{Inhaltsverzeichnis}
\tableofcontents

% arabische Seitenzahlen im Hauptteil ------------------------------------------------------------------
\clearpage{}
\pagenumbering{arabic}
\setcounter{page}{2} %%% Dieser Pagecounter muss entsprechend der verbrauchten Seiten im Inhaltsverzeichnis angepasst werden. Endet das IHV bei Seite III, so muss hier 4 eingetragen werden



% Inhalt -----------------------------------------------------------------------------------------------
% ------------------------------------------------------------------------------------------------------

% Kapitel sinnvollerweise in eigene Dateien schreiben und einbinden
\section{Einleitung}
Lorem ipsum dolor, sit amet consectetur adipisicing elit.
Quis deleniti, velit distinctio qui inventore id aliquid illum,
sunt perspiciatis unde dicta iusto, officiis eius eos fuga possimus sit omnis ratione?
% Neue Seite nach Kapitel
\clearpage{}



% Literaturverzeichnis ---------------------------------------------------------------------------------
%   Das Literaturverzeichnis wird ggf. aus der BibTeX-Datenbank "document.bib" erstellt.
% ------------------------------------------------------------------------------------------------------
% Quelle ins Quellenverzeichnis aufnehmen, die im Dokument nicht genutzt wird.
\nocite{logo_hska}
\renewcommand\refname{Quellen}
\printbibliography



% Anhang -----------------------------------------------------------------------------------------------
%\begin{appendix}
 %   \clearpage{}
  %  \pagenumbering{roman}
%	\section{Anhang}
%\end{appendix}

\end{document}